\mychapter{Conclusões}{cap:conclusion}\lhead{CONCLUSÕES}

Esta dissertação investiga se características da demanda ajudam a escolher políticas de reposição com melhor compromisso entre custo e atendimento. O AIPE organiza essa avaliação em preparação, caracterização e comparação de candidatos. Os resultados disponíveis sustentam examinar a utilidade incremental da escolha por perfil; ainda não demonstram economia geral nem validam o PSE supervisionado.

\section{Respostas às questões de pesquisa}

O desempenho relativo das políticas varia entre recortes e períodos. O resultado favorável de GA--DQN no piloto da Paraíba não se repete como superioridade agregada na Bahia. No M5 histórico, apenas SA e PPO superam o serviço médio de 0,70. Essas diferenças descrevem implementações e protocolos específicos, sem identificar um controlador universal.

A nova execução da Bahia torna mais precisa a resposta sobre os perfis. Os 4.350 ajustes cobrem as 145 séries em duas origens, com escolha na validação das demais lojas e avaliação no teste da loja reservada. Em conjunto, selecionar por perfil aumenta o custo em 1,93\% e o atendimento médio em 3,89 pontos percentuais. Os casos abaixo da meta passam de 119 para 100 entre os 283 com demanda positiva no teste. O resultado representa serviço adicional com custo maior. Na segunda origem, as estratégias fazem escolhas idênticas.

Essas observações ainda não respondem integralmente às perguntas sobre aprendizagem e generalização. A escolha explícita por perfil foi avaliada temporalmente, mas o PSE supervisionado requer treinamento e teste próprios. A coorte continua selecionada retrospectivamente; as duas janelas são curtas, compartilham um ciclo e repetem as mesmas séries. A transferência para outros produtos, períodos e fontes de dados permanece aberta.

\section{Contribuições e limites}

A contribuição metodológica é distinguir três problemas: calibrar uma política, escolher uma candidata do portfólio e avaliar a recomendação em dados posteriores. A contribuição empírica é examinar essa distinção em dados reais do varejo brasileiro e contrastá-la com resultados históricos públicos do M5. A infraestrutura integrada permite rastrear dados, configurações, sementes e saídas da nova execução.

Os métodos do portfólio e a arquitetura GA--DQN já existem na literatura. Seu uso não constitui novidade algorítmica. Na nova Bahia, o híbrido melhora custo e atendimento frente ao DQN, mas não frente ao GA na média dos testes. Como também consome busca genética e demonstrações, a comparação não isola um ganho causal da combinação. A avaliação temporal M5 deverá complementar essa leitura antes de uma decisão definitiva sobre o portfólio.

As matrizes históricas preservam valor exploratório, com limites de rastreabilidade, uso do histórico completo e repetições determinísticas. A nova execução restringe o ajuste ao passado e documenta sementes independentes de treinamento, mas essas sementes não são novas lojas nem novas realizações da demanda. Custos simulados, vendas possivelmente censuradas e estados iniciais de referência delimitam a interpretação operacional.

\section{Implicações e continuidade}

A definição do serviço antecede a escolha do controlador. Média por série, atendimento ponderado pela demanda e garantia individual produzem exigências distintas. Na nova Bahia, as médias das estratégias superam 0,70 quando as origens são reunidas, mas o atendimento ponderado permanece abaixo desse alvo e há numerosas violações individuais. A média global, portanto, é insuficiente para julgar a adequação empresarial.

O avanço seguinte é consolidar o M5 sob o protocolo temporal, examinar estabilidade temporal e definir o valor econômico do atendimento adicional antes de treinar o PSE. O plano do Capítulo~\ref{cap:plano_continuidade} organiza essas dependências. A conclusão atual é delimitada: perfis podem alterar o compromisso entre custo e serviço, mas seu benefício precisa ser demonstrado frente a uma política global competitiva em períodos reservados.
